\documentclass[11pt,a4paper]{article}

\usepackage[margin=2.2cm]{geometry}
\usepackage{graphicx}
\usepackage{booktabs}
\usepackage{amsmath,amssymb}
\usepackage[font=small,labelfont=bf]{caption}
\usepackage{xcolor}
\usepackage[colorlinks=true,linkcolor=blue!60!black,urlcolor=blue!60!black,citecolor=blue!60!black]{hyperref}
\usepackage{listings}
\usepackage{enumitem}
\usepackage{fancyhdr}

\setlength{\parskip}{4pt}
\setlength{\parindent}{0pt}
\graphicspath{{figures/}}

\definecolor{tabhdr}{HTML}{E8E8E8}
\definecolor{wingreen}{HTML}{2E7D32}
\definecolor{winred}{HTML}{C62828}

\pagestyle{fancy}
\fancyhf{}
\fancyhead[L]{\small SolarFlare V4 --- S-series ablations}
\fancyhead[R]{\small harp\_11930 focus}
\fancyfoot[C]{\thepage}
\renewcommand{\headrulewidth}{0.4pt}

\title{\textbf{SolarFlare V4 SimpleConvLSTM:\\ S-series Ablations Report}\\
       \large Predicting Winding-Flux Extremes from Sequential Magnetograms\\
       \large Case Study: HARP 11930}
\author{Manik Sharma\\\small{Engineering Physics, IIT BHU}}
\date{\today}

\begin{document}
\maketitle
\thispagestyle{empty}

\begin{abstract}
\noindent
We report an ablation sweep over a minimal 2-layer convolutional LSTM
(\emph{SimpleConvLSTM}) for forecasting the next 4 frames of winding
flux \(w(\mathbf{x},t)\) of an active region. The goal of the
sweep is to climb above the trivial \emph{persistence baseline}
(``predict the last frame'') on the Critical Success Index (CSI) of
extreme flux pixels --- a thresholded binary metric directly relevant
to flare onset.
We trained 16 sister models (S0--S16, excluding two that re-ran already-known
configurations). The first eleven (S0--S11) varied loss
re\-weighting, normalisation, teacher-forcing schedule, and field
representation; \textbf{all eleven lost to persistence}. We then pivoted
to a \emph{dual-head architecture} that adds a per-pixel binary
classifier (``is this pixel extreme at time \(t\)?'') trained with
class-imbalance-corrected binary cross-entropy alongside the regression
loss. Arm \textbf{S13} first matched persistence
(CSI~\(\sim\)0.043) and \textbf{S16} --- the hybrid of S13's classifier
head with S11's extreme-pixel-weighted regression --- \textbf{beat
persistence by \(\sim\)31\%} (CSI 0.0562 vs.\ persistence 0.043) while
simultaneously improving regression CSI 2.8\(\times\) over S13. All
metrics are reported on a fixed informative test set of three flare-rich
HARP cubes [245, 274, 49]. We close with a multi-step autoregressive
``staircase'' diagnostic on HARP 11930 showing that even the best
single-step arms either mode-collapse (S3) or runaway-explode (S11)
under chained prediction, indicating regression fidelity at horizons
longer than 4 frames remains open.
\end{abstract}

\tableofcontents
\newpage

\section{Notation \& common setup}

\textbf{Task.} Given last \(T_{\text{in}}=10\) frames of winding flux
\(w(\mathbf{x},t)\) over a SHARP cube, predict next \(T_{\text{out}}=4\)
frames (48-min lead). All arms use the same 2-layer ConvLSTM
(\texttt{SimpleConvLSTM}, kernel 3, hidden 64, residual decode
\(\hat w_t = w_{t-1} + \Delta\hat w_t\), \(\sim\)0.9\,M params).

\textbf{Data.} 21 SHARP cubes, 12-min cadence; \(64\!\times\!64\) crops
stride 4; per-cube z-score normalisation (S7 alternative: signed
\(\operatorname{asinh}\)); fixed informative test set
\{HARP 245, 274, 49\} from S8 onwards; 0.90/0.05/0.05
train/val/test split.

\textbf{Training.} Adam, batch 8, \(\eta=10^{-3}\); 15 epochs;
patience 8 on val loss; teacher-forcing fraction
\(\tau\!:\!1\!\to\!0\) over first \texttt{tf\_decay\_epochs} epochs;
gradient clip 0.5; cosine LR schedule (S0--S15) or constant (S16+);
no augmentation in dual-head arms.

\textbf{Metrics on extreme mask \(|w|>0.528\sigma\) (99th-pct):}
\begin{itemize}[leftmargin=14pt,itemsep=1pt,topsep=2pt]
\item \textbf{CSI} \(=\text{TP}/(\text{TP}+\text{FP}+\text{FN})\), our
      primary skill score.
\item \textbf{CLS-CSI} = CSI on the classifier-head sigmoid
      \(>\!0.5\) (dual-head arms only).
\item \textbf{Persistence baseline} predicts the most recent input
      frame for all 4 future steps; achieves
      \textbf{CSI 0.043, HSS 0.082} on this test set ---
      \emph{the wall to beat}.
\item Auxiliary: SSIM (\texttt{data\_range=2}), persistence skill
      \(=(1-\text{MAE}_{\text{mod}}/\text{MAE}_{\text{pers}})\!\times\!100\%\),
      variation ratio
      \(\text{mean}|\Delta\hat w|/\text{mean}|\Delta w|\) (mode collapse if
      \(\to 0\), runaway if \(\gg 1\)).
\end{itemize}

\textbf{Losses used across arms.}
\begin{itemize}[leftmargin=14pt,itemsep=1pt,topsep=2pt]
\item \(\mathcal{L}_{\text{L1}} = \langle |\hat w - w| \rangle\)
\item \(\mathcal{L}_{\text{wMAE}} = \langle m\!\cdot\!|\hat w - w| \rangle\),
      \(m{=}\kappa\) where \(|w|>\tau\), else 1
      (\(\kappa{=}\)\texttt{extreme\_pixel\_weight}).
\item \(\mathcal{L}_{\text{BCE}} =
      \text{BCE}(\ell, \mathbb{1}_{|w|>\tau};
      \texttt{pos\_weight}{=}\alpha)\) on classifier logits.
\item \(\mathcal{L}_{\text{composite}}\) (S3 only):
      L1 + SSIM-loss + AsymmetricExtreme + temporal-difference
      + temporal-variance + extreme-pixel weighting.
\item Dual-head total:
      \(\alpha_r\mathcal{L}_{\text{(w)MAE}} + \beta\mathcal{L}_{\text{BCE}}\).
\end{itemize}

\newpage
\section{Cross-arm parameter matrix}

The S-series is a single-axis-at-a-time sweep starting from the
minimal baseline S0. The table below is the chronological list of arms
included in this report; each row notes only the delta against the
preceding arm.

\begin{table}[h]
\centering\footnotesize
\begin{tabular}{@{}l p{5.4cm} p{8.2cm}@{}}
\toprule
Arm & Delta vs.\ prior & Question being asked \\
\midrule
\rowcolor{tabhdr}S0 & baseline (L1, D4 aug)              & Does the 2-layer ConvLSTM by itself match the 6-layer attention model? \\
S2 & + residual decode                                  & Does \(w_t = w_{t-1} + \Delta\) anchor predictions to persistence? \\
\rowcolor{tabhdr}S4 & S2 + teacher-forcing (\(\tau{:}1\!\to\!0\)) & Does TF cure autoregressive explosion of S2? \\
S5 & S4 + 90/5/5 split + GroupNorm                       & Does more data lift CSI? \\
\rowcolor{tabhdr}S3 & S5 + 6-term composite loss          & Does kitchen-sink loss break L1 mean-collapse? \\
S6 & S5 + \texttt{extreme\_pixel\_weight}{=}25 only      & Which composite term drove S3? \\
\rowcolor{tabhdr}S8 & S5 on fixed test cubes (L1)         & Recalibrate against persistence on a fixed evaluation set. \\
S7 & S8 + signed-asinh normalisation                    & Does heavy-tail compression help CSI? \\
\rowcolor{tabhdr}S9 & S8 + BatchNorm (was S4's setting)   & Is BatchNorm the source of S4 ``sharpness''? \\
S10 & S8 + \texttt{tf\_decay\_epochs}{=}5 + patience 8    & Does a complete TF curriculum help? \\
\rowcolor{tabhdr}S11 & S10 + extreme\_pixel\_weight 25    & Both working levers stacked --- does the combo beat persistence? \\
\midrule
\multicolumn{3}{l}{\textit{Dual-head architecture pivot (classifier head added; constant LR from S16).}}\\
\midrule
\rowcolor{tabhdr}S13 & dual-head, \(\alpha{=}1, \beta{=}1\), \texttt{pos\_weight}{=}100 & Does an explicit per-pixel BCE classifier escape regression mean-collapse? \\
S14 & S13 + focal loss (\(\gamma{=}2, \alpha_f{=}0.25\)) & Does focal weighting outperform pos\_weight BCE? \\
\rowcolor{tabhdr}S15 & S13 + \(\alpha{=}0.1\) (regression-dominated $\to$ classifier-dominated) & Does starving the regression head free capacity for the classifier? \\
S16 & S13 with \texttt{pos\_weight}{=}60 + \texttt{extreme\_pixel\_weight}{=}25 & Joint train of S11 regression lever \(+\) S13 classifier head. \\
\bottomrule
\end{tabular}
\caption{Chronological sweep. S1 (no-aug), S12 (cosine-LR dual-head),
S17 (post-hoc inference combo, no training) discussed where relevant
in the text. S18+ excluded from this report.}
\end{table}

All arms terminated after 15 epochs (or earlier on patience). All
metrics reported below are on the fixed test set
\{HARP 245, 274, 49\} unless otherwise noted; figures show
HARP 11930 (\emph{train} cube, qualitative only).

\newpage
\section{Regression-only arms (S0--S11)}

Eleven single-head arms varying loss, normalisation, teacher-forcing
curriculum, and field representation. Headline: every arm loses to
persistence (CSI 0.043). Best regression CSI = 0.0176 (S11) ---
41\% of persistence.

\subsection{S0 --- minimal baseline}
\textbf{Loss:} L1 only. \textbf{Aug:} D4. \textbf{Norm:} BatchNorm.
\textbf{Test CSI:} 0 across all epochs; best val 0.039 at ep4 then
overfit.\\
\textbf{Finding:} The 2-layer ConvLSTM with plain L1 collapses to the
zero-field conditional median. Pred amplitude never crosses the
extreme threshold so TP=0. Confirmed the wall is structural, not a
capacity issue --- the deep 6-layer model also gets CSI \(\sim\)0.01.
\\[2pt]\noindent\includegraphics[width=\linewidth]{s0_harp11930_fullframe}\\[2pt]
\noindent\includegraphics[width=\linewidth]{S0_loss}

\subsection{S2 --- residual decode}
\textbf{Delta:} predict \(\Delta w_t\), output \(w_{t-1}+\Delta\).
\textbf{Finding:} Without teacher forcing this \emph{explodes} ---
variance ratio climbs 18\(\to\)764 across 15 epochs (autoregressive
positive feedback). Confirmed residual anchor needs a stabiliser.
(No figure: model output unusable.)

\subsection{S4 --- S2 + teacher forcing}
\textbf{Delta:} TF \(\tau{=}1{\to}0\) over training.
\textbf{Test:} CSI 0.099, HSS 0.178, SSIM 0.438 (old random split).
\textbf{Finding:} TF cures the explosion; variance ratio settles near
1. Looks ``sharp'' visually but on the fixed test set later this
sharpness turns out to be uncontrolled BatchNorm variance, not skill
(see S9). \\[2pt]\noindent\includegraphics[width=\linewidth]{s4_harp11930_fullframe}

\subsection{S5 --- S4 + 90/5/5 split + GroupNorm}
\textbf{Delta:} more training data, swap BN\(\to\)GN.
\textbf{Test:} SSIM 0.70, persistence skill +30\%, CSI \(\approx\) 0.
\textbf{Finding:} GroupNorm + more data improves spatial similarity
and beats persistence MAE, but \emph{smoother} field $\Rightarrow$
\emph{lower} CSI on extremes. Trade-off explicit: L1 rewards smoothing,
extreme classification punishes it.\\[2pt]\noindent\includegraphics[width=\linewidth]{s5_harp11930_fullframe}\\[2pt]
\noindent\includegraphics[width=\linewidth]{S5_loss}

\subsection{S3 --- S5 + 6-term composite loss}
\textbf{Loss:} L1 + SSIM + AsymExtreme + TempDiff + TempVar + extreme-pixel weight 25.
\textbf{Test:} best val CSI 0.0236 (ep13); test 0 (random split put a
flareless cube in test). \textbf{Finding:} Composite is the best
\emph{regression} on the qualitative HARP 11930 spatial integral (see
\S\ref{sec:cmp_intflux}) but the lift over single-term arms is small
($\sim$2\(\times\) over S0). S6 below isolates which term carries the
signal.\\[2pt]\noindent\includegraphics[width=\linewidth]{s3_harp11930_fullframe}\\[2pt]
\noindent\includegraphics[width=\linewidth]{S3_loss}

\subsection{S6 --- only the extreme-pixel weight term}
\textbf{Loss:} L1 + \texttt{extreme\_pixel\_weight}{=}25 (everything else off).
\textbf{Test:} val CSI 0.0215 \(\approx\) S3.
\textbf{Finding:} S3's full composite is \emph{not} additive --- 95\%
of its lift comes from the single extreme-pixel-weight term. SSIM /
asymmetric / temporal terms add \(\sim\)0. Identifies the only loss
lever that moves the needle.\\[2pt]\noindent\includegraphics[width=\linewidth]{s6_harp11930_fullframe}\\[2pt]
\noindent\includegraphics[width=\linewidth]{S6_loss}

\subsection{S8 --- S5 (L1) on fixed test cubes}
\textbf{Delta:} pin test set to \{245, 274, 49\}.
\textbf{Test CSI:} 0.0033 vs persistence 0.043 (\(13\times\) below).
\textbf{Finding:} Establishes the persistence wall on a flare-rich
fixed test set; the optimism in S4's old random-split CSI 0.099 was
a flareless test-cube artefact.\\[2pt]\noindent\includegraphics[width=\linewidth]{s8_harp11930_fullframe}\\[2pt]
\noindent\includegraphics[width=\linewidth]{S8_loss}

\subsection{S7 --- S8 + signed-asinh normalisation}
\textbf{Delta:} swap z-score for
\(\operatorname{sign}(w)\operatorname{asinh}(|w|/s)\) representation
(\textbf{not} the loss; only how data enters the model). \textbf{Test
CSI:} 0.0042 \(\approx\) S8. \textbf{Finding:} Heavy-tail input
compression has no effect on CSI. Apparent SSIM 0.98 is an
asinh-space artefact, not improvement in physical space. Representation
lever \textbf{rejected}.
\\[2pt]\noindent\includegraphics[width=\linewidth]{s7_harp11930_fullframe}\\[2pt]
\noindent\includegraphics[width=\linewidth]{S7_loss}

\subsection{S9 --- S8 with BatchNorm (= S4 on fixed test)}
\textbf{Delta:} swap GN\(\to\)BN.
\textbf{Test CSI:} 0.0118 (\(3.5\times\) S8 \emph{but} still below
persistence by 4\(\times\)). \textbf{Finding:} BN drives over-variance
(ratio 3.7, val hit 19); unstable (SSIM 0.51\(\to\)0.09, early-stop
ep5). S4's visual ``sharpness'' was uncontrolled variance, \textbf{not}
forecast skill. BN rejected.\\[2pt]\noindent\includegraphics[width=\linewidth]{s9_harp11930_fullframe}\\[2pt]
\noindent\includegraphics[width=\linewidth]{S9_loss}

\subsection{S10 --- S8 + fast TF curriculum}
\textbf{Delta:} \texttt{tf\_decay\_epochs}{=}5, patience 8.
\textbf{Test CSI:} 0.0027 \(\approx\) S8.
\textbf{Finding:} Free-running TF curriculum bumps variance ratio to
0.75 briefly at ep7 but L1 drags it back to 0.028 by ep15. TF is
necessary-but-insufficient on its own; L1 mean-collapse defeats the
de-smoothing.\\[2pt]\noindent\includegraphics[width=\linewidth]{s10_harp11930_fullframe}\\[2pt]
\noindent\includegraphics[width=\linewidth]{S10_loss}

\subsection{S11 --- S10 fast-TF \(+\) extreme weight (both working levers)}
\textbf{Delta:} S10 + S6's extreme\_pixel\_weight 25.
\textbf{Test CSI:} \textbf{0.0176} --- best of the S0--S11 series
(5\(\times\) S8, 41\% of persistence). \textbf{Finding:} The two
working levers are additive in CSI but the combo \emph{still} loses to
persistence by 2.4\(\times\). Variance ratio re-collapses to 0.036 by
ep15 --- L1 wins long-run.\\[2pt]\noindent\includegraphics[width=\linewidth]{s11_harp11930_fullframe}\\[2pt]
\noindent\includegraphics[width=\linewidth]{S11_loss}
\label{fig:s11_loss}

\subsection*{Verdict for the regression-only family}
Four levers explored (loss reweighting, BN/GN, TF curriculum,
input representation). The only loss term that moves CSI is
\texttt{extreme\_pixel\_weight}; the only TF schedule that helps is
the full \(\tau{:}1\!\to\!0\). Stacking both (S11) gets to 41\% of
persistence. \textbf{All single-head paths exhausted}; the model
cannot push regression amplitude past the extreme threshold without
also losing the L1 fit elsewhere.
\textbf{Conclusion: structural pivot needed.}

\newpage
\section{Dual-head pivot (S13--S16)}\label{sec:s13}

Pivot rationale: the extreme-mask CSI is a thresholded binary, but
S0--S11 train a regression loss on amplitude. Mismatch fixed by
adding a parallel per-pixel BCE classifier head on the same ConvLSTM
trunk:
\(
\hat w(\mathbf{x},t),\ \ell(\mathbf{x},t)
= \text{ConvLSTM}(\mathbf{w}_{1:T_{\text{in}}})
\),
with training loss
\(
\mathcal{L} = \alpha\,\mathcal{L}_{\text{(w)MAE}}(\hat w, w)
            + \beta\,\mathcal{L}_{\text{BCE}}(\ell, \mathbb{1}_{|w|>\tau}; \texttt{pos\_weight}{=}p)
\).
At test time we report both \textbf{regression CSI} (CSI on
\(\hat w\)) and \textbf{CLS-CSI} (CSI on \(\sigma(\ell)>0.5\)).
\textbf{Constant LR \(\eta=10^{-3}\)} (no cosine) from S16 onwards
after observing cosine washes out late-epoch gradient signal at our
short 15-epoch schedule.

\textit{S12 (cosine LR rerun) and S15 (classifier-dominant
\(\alpha{=}0.1\)) below either confirmed expected behaviour or did not
finish on the planned slot and are summarised in text only.}

\subsection{S13 --- first to match persistence}
\textbf{Loss:} \(\alpha{=}1,\beta{=}1,p{=}100\), plain L1 on
regression head. \textbf{LR:} constant
\(10^{-3}\). \textbf{15 epochs.}\\
\textbf{Test:} \textbf{CLS-CSI 0.0434, HSS 0.0749} (matches
persistence 0.043); regression CSI 0.0090 (dead --- \(p{=}100\)
starves L1 budget). Persistence skill +35.7\%. Val loss
0.164\(\to\)0.149 monotone-ish; val CLS-CSI 0.003\(\to\)0.007.\\
\textbf{Finding:} \emph{First} S-arm to match persistence on the
extreme mask. The classifier head is the load-bearing
component; the regression head exists mainly to provide a spatial
backbone for the classifier features.
\\[2pt]\noindent\includegraphics[width=\linewidth]{s13_harp11930_fullframe}

\subsection{S14 --- focal loss is the wrong instrument}
\textbf{Loss:} same as S13 but BCE \(\to\) focal
\(\bigl(-\alpha_f(1-p)^\gamma\log p\bigr)\),
\(\gamma{=}2,\alpha_f{=}0.25\), \(p{=}60\).\\
\textbf{Test:} \textbf{CLS-CSI 0 every epoch.} Best val 0.0072 ep10
(focal numeric scale is just small, not a learning signal). Regression
SSIM 0.79--0.91, same as S0--S11.\\
\textbf{Finding:} Focal's \((1-p)^\gamma\) factor over-suppresses
gradients on the rare positives at our \(\sim\)0.1\% imbalance level.
The classifier never gets enough signal to leave its initialisation.
\textbf{Focal rejected for this task.} (Figure omitted --- the
regression head produces the same pattern as S0--S11; the classifier is
dead.)

\subsection{S15 --- classifier-dominant (\(\alpha=0.1\))}
\textbf{Loss:} \(\alpha{=}0.1, \beta{=}1, p{=}60\). Idea: starve the
regression head to free capacity for the classifier.\\
\textbf{Test:} CLS-CSI 0.0445 (\(\approx\) S13's 0.0434, no lift);
regression CSI 0.0184 (above S13's 0.0090 --- the small regression
weight still produces a useful head, somewhat counter-intuitively).\\
\textbf{Finding:} Lowering \(\alpha\) does \emph{not} buy classifier
gain. Confirms classifier head saturates above some \(\beta\)
threshold and the remaining headroom is on the regression side, not
classifier capacity. Recipe pointer for the winning hybrid (S16).

\subsection{S16 --- joint hybrid: weighted-MAE + classifier}\label{sec:s16}
\textbf{Loss:} \(\alpha{=}1, \beta{=}1, p{=}60\),
\(\mathcal{L}_{\text{wMAE}}\) with
\texttt{extreme\_pixel\_weight}{=}25 on the regression head
(S11's lever) plus BCE classifier (S13's lever). Code change:
extended \texttt{training/losses/dual\_head.py} to take an
\texttt{extreme\_pixel\_weight} kwarg and swap
\texttt{F.l1\_loss} for \texttt{WeightedMAELoss} when \(>1\).\\

\textbf{Test:}
\begin{center}\small
\begin{tabular}{lccc}
\toprule
 & reg.\ CSI & \textbf{CLS-CSI} & cls HSS \\
\midrule
persistence  & ---     & 0.043 & 0.082 \\
S11 (best regression)  & 0.0176 & ---    & ---    \\
S13 (best classifier alone)  & 0.0090 & 0.0434 & 0.0749 \\
\textbf{S16 hybrid}  & \textbf{0.0255} & \textbf{0.0562} & \textbf{0.0995}\\
\bottomrule
\end{tabular}
\end{center}

\textbf{Finding:}
S16 is the first S-arm to \emph{beat} persistence on the extreme mask
(\(\sim\)31\% above; not just match it like S13). Regression CSI is
2.8\(\times\) S13 and 1.4\(\times\) S11. The extreme-pixel-weighted MAE
rescues the regression head's amplitude (briefly crosses the threshold
ep8--9 producing nonzero reg CSI in-training) while the lighter
\(p{=}60\) leaves enough budget for the BCE classifier to learn.
Persistence skill 34\%, SSIM 0.69 --- comparable to S13 on auxiliary
metrics. Hybrid hypothesis confirmed: the two heads' levers compose,
they do not antagonise.
\\[2pt]\noindent\includegraphics[width=\linewidth]{s16_harp11930_fullframe}

\newpage
\section{S17 --- post-hoc S11\(\times\)S13 inference combo}

Before committing GPU to S16 training we asked the cheaper
question: \emph{would} multiplying S11's regression output by S13's
classifier confidence improve the joint forecast? No training; both
checkpoints already on disk. Three fusion variants tested on HARP
11930, full-frame tiled (3 active windows, integrated winding-flux
MAE as proxy for end-to-end skill):

\begin{center}\small
\begin{tabular}{l c c}
\toprule
Variant & Mean MAE & vs.\ S11 alone \\
\midrule
\rowcolor{tabhdr}
\textbf{B\_soft:} \(\hat w_{\text{S11}} \cdot \bigl(1 + 2\sigma(\ell_{\text{S13}})\bigr)\)
& \textbf{2.06\(\times\)10\(^8\)} & \textcolor{wingreen}{\textbf{\(-\)8\%}}\\
S11 alone (regression)             & 2.23\(\times\)10\(^8\) & --- \\
A\_gate: \(\hat w_{\text{S11}} \cdot (\sigma(\ell_{\text{S13}}) > 0.5)\)
& 2.38\(\times\)10\(^8\) & \textcolor{winred}{+7\%} \\
S13 alone (regression head only) & 3.00\(\times\)10\(^8\) & +34\% \\
C\_blend: S13 where mask else S11 & 3.01\(\times\)10\(^8\) & +35\% \\
Persistence (pred = 0)             & 4.84\(\times\)10\(^8\) & +117\% \\
\bottomrule
\end{tabular}
\end{center}

Mask coverage 5.5--5.9\% of pixels per window (classifier fires on a
small fraction). Lessons:
\begin{itemize}[leftmargin=14pt,itemsep=1pt,topsep=2pt]
\item \textbf{Soft boost \(>\) hard gate.} Zeroing background flux
      hurts integrated MAE more than amplifying extreme zones helps.
\item \textbf{S13 regression alone is useless} (pos\_weight 100 left
      no L1 budget); blending S13 in for extreme zones \emph{worsens}
      the integrated flux relative to S11 alone.
\item \textbf{Hybrid hypothesis confirmed} at the post-hoc level
      (8\% lift). The natural next question --- ``does joint
      training recover this gain at the loss level?'' --- is exactly
      what S16 (\S\ref{sec:s16}) answers: yes, and it
      \emph{compounds} (S16 classifier CSI 0.0562 = 29\% over S13's
      0.0434, alongside regression CSI 2.8\(\times\) better).
\end{itemize}

\newpage
\section{Multi-step autoregressive ``staircase'' diagnostic}\label{sec:cmp_intflux}

Single-step CSI evaluates the model on the next 4 frames only.
For mission-relevant lead times we need the model to chain its own
predictions. We instantiated a staircase protocol on HARP 11930
(\emph{train} cube, qualitative): input \(T_{\text{in}}=10\) real
frames, predict \(T_{\text{out}}=4\), commit the first
\texttt{stride}\(=2\) predicted frames into the input buffer, slide,
repeat for 20 steps (i.e.\ 40 committed frames \(\equiv\) 8\,h forecast
lead). Full-frame inference at each step (tiled, stride 64, no
overlap). Pure autoregressive --- the model never sees a real frame
past \(T_{\text{in}}\) after step 0.

We ran the protocol on S3 (best single-step regression on this cube),
S11 (best single-step regression on the test cubes), and S16 (new
SOTA).

\subsection*{Headline numbers}
\begin{center}\small
\begin{tabular}{l r r r}
\toprule
Arm & MAE (integrated flux) & Var.\ ratio & Drift corr.\ \(\rho\)\\
\midrule
S3  & \(1.5\!\times\!10^8\) & 0.028 (mode collapse)      & +0.14 \\
S11 & \(7.6\!\times\!10^8\) & 15.8 (runaway explosion)   & +0.81 \\
S16 & (see below)                 & ---                       & ---   \\
\bottomrule
\end{tabular}
\end{center}

\emph{(S16 staircase numbers are in \texttt{outputs/staircase\_s16\_harp\_11930/series.npz}; integrated-flux plot below.)}

\subsection*{S3: mode collapse}
Within 4 chained frames the predicted field amplitude drops from
\(|\hat w|_\infty\!=\!1.3\!\times\!10^7\) to
\(\sim\!10^5\) (100\(\times\) collapse). Frame-to-frame field
difference settles to \(\sim\)150 on values of order \(1.5\!\times\!10^5\)
(\(<\)0.1\% relative motion). After collapse the network produces a
\emph{static} low-amplitude field for the remaining 36 frames. MAE
stays low only because the static value happens to be close to the
GT temporal mean.

\noindent\includegraphics[width=0.49\linewidth]{s3_staircase_4steps}
\includegraphics[width=0.49\linewidth]{s3_staircase_intflux}

\subsection*{S11: runaway explosion}
Behaves like S3 for the first \(\sim\)9 chained steps (low-amplitude
collapse phase) and then \emph{variance explodes}: frame
\(|\hat w|_\infty\) grows from \(\sim\!10^5\) to
\(2\!\times\!10^6\) by frame 39; per-frame field difference grows from
150 to \(\sim\)3000. The 12-step grid below makes the transition
visible --- the bottom rows are progressively more saturated.
Variance ratio 15.8\(\times\), drift correlation +0.81 (error grows
strongly with horizon).

\noindent\includegraphics[width=0.49\linewidth]{s11_staircase_4steps}
\includegraphics[width=0.49\linewidth]{s11_staircase_intflux}

\noindent\includegraphics[width=\linewidth]{s11_staircase_12steps}

\subsection*{S16: staircase}
\noindent\includegraphics[width=0.49\linewidth]{s16_staircase_4steps}
\includegraphics[width=0.49\linewidth]{s16_staircase_intflux}

\noindent\includegraphics[width=\linewidth]{s16_staircase_12steps}

\subsection*{Verdict}
None of S3 / S11 sustains chained prediction --- one collapses, the
other explodes. The two failure modes are symmetric: L1-dominated
losses converge to a stable low-amplitude fixed point, while
amplitude-rewarding losses (S11's extreme weight + fast-TF) shift
the autoregressive map to an unstable regime. Long-horizon
forecasting is genuinely open; the single-step gains from S16 should
be re-tested under chained rollouts before claiming
multi-hour skill.

\newpage
\section{Cross-arm summary \& next levers}

\subsection*{Test-set CSI ranked (fixed test cubes 245, 274, 49)}
\begin{center}\small
\begin{tabular}{l r r r l}
\toprule
Arm & reg.\ CSI & CLS-CSI & cls HSS & vs.\ pers.\ 0.043 \\
\midrule
\rowcolor{tabhdr}\textbf{S16} (hybrid wMAE + classifier)  & \textbf{0.0255} & \textbf{0.0562} & \textbf{0.0995} & \textcolor{wingreen}{\textbf{+31\% above}} \\
S15 (\(\alpha{=}0.1\) dual)               & 0.0184 & 0.0445 & 0.0783 & matches \\
\rowcolor{tabhdr}S13 (dual, \(p{=}100\))   & 0.0090 & 0.0434 & 0.0749 & matches \\
S11 (regression only, fastTF + ext.\ wt.) & 0.0176 & --- & --- & 41\% of \\
S6 (extreme weight only)                  & 0.0215 (val) & --- & --- & --- \\
S3 (full composite)                       & 0.0236 (val) & --- & --- & --- \\
S9 (BatchNorm)                            & 0.0118 & --- & --- & 27\% of \\
S7 (signed asinh)                         & 0.0042 & --- & --- & 10\% of \\
S8 (L1 fixed test)                        & 0.0033 & --- & --- & 8\% of \\
S10 (fast-TF L1)                          & 0.0027 & --- & --- & 6\% of \\
S0 (baseline)                             & 0 & --- & --- & 0\% \\
S14 (focal classifier)                    & 0 reg, \textbf{0 cls} & 0 & 0 & dead arm \\
\bottomrule
\end{tabular}
\end{center}

\subsection*{Mechanism narrative (one sentence per arm cluster)}
\begin{itemize}[leftmargin=14pt,itemsep=2pt,topsep=2pt]
\item \textbf{S0--S11 (regression-only)}: L1 minimises the conditional
      median of a heavy-tailed field; predictions never cross the
      extreme threshold; CSI capped well below persistence.
\item \textbf{S13 (dual-head, pos\_w 100)}: explicit BCE classifier
      bypasses the regression bottleneck; matches persistence at the
      cost of a dead regression head.
\item \textbf{S14 (focal)}: focal's \((1-p)^\gamma\) factor over-suppresses
      gradients on rare positives at our imbalance level $\Rightarrow$
      classifier never escapes init.
\item \textbf{S15 (\(\alpha{=}0.1\))}: lowering regression weight does not
      buy classifier headroom; classifier saturates above some
      \(\beta\)-threshold.
\item \textbf{S16 (joint hybrid)}: weighted-MAE on regression head
      + lighter pos\_w 60 BCE classifier $\Rightarrow$ both heads
      learn complementary signals; first arm above persistence.
\end{itemize}

\subsection*{Next levers}
If we want to push past S16's +31\%-of-persistence:
\begin{enumerate}[leftmargin=18pt,itemsep=2pt,topsep=2pt]
\item \texttt{pos\_weight} sweep at fixed
      \texttt{extreme\_pixel\_weight}{=}25: find the Pareto
      frontier between regression and classifier heads.
\item \textbf{Quantile regression head} (pinball loss at
      \(\tau\in\{0.5,0.9,0.99\}\)): predict the 99-th percentile
      amplitude directly rather than gating an L1 mean. Skips the
      classifier detour.
\item \textbf{Event-onset reframe}: per-frame
      \(P(\exists\,\text{extreme pixel in next K frames})\) using
      GOES-aligned event labels --- orthogonal to per-pixel rate; uses
      the half of HARP cubes (e.g.\ 11930) where pixel rate and GOES
      events are decoupled.
\item \textbf{GradNorm balancing}: auto-balance
      \(\alpha, \beta\) so per-batch gradient norms are equal; avoids
      hand tuning the multi-task loss weights.
\item \textbf{Autoregressive regularisation}: penalise variance ratio
      \(\to\)1 during training, addressing both S3 collapse and S11
      explosion identified in the staircase diagnostic (\S\ref{sec:cmp_intflux}).
\end{enumerate}

\subsection*{Reproducibility pointers}
Configs: \texttt{configs/ablations/S\{0..16\}\_*.yaml}.
Best checkpoints: \texttt{outputs/ablations/<arm>/checkpoints/best\_model.pt}.
Compare \& staircase scripts:
\texttt{scripts/s0\_viz/\{\_compare\_harp11930, \_s17\_combo,
\_staircase\_harp11930, \_staircase\_viz\}.py}. Per-epoch logs:
\texttt{logs/main\_\_5060ti\_cuda.log} and
\texttt{logs/main\_\_mini\_mps.log}.

Full per-experiment writeup
in \texttt{docs/V4/findings\_dual\_head.md} (dual-head pivot),
\texttt{docs/V4/findings\_simple\_convlstm\_S0\_S1.md} (early regression
arms), and the ablation matrix table in
\texttt{docs/V4/04\_experiments\_and\_state.md}.


\end{document}
